\section{Introduction}
Iterative Bregman projections turn structured constrained optimization into alternating normalization steps. Their sparse linear algebra and GPU-compatible updates are useful in machine learning, notably for optimal transport, barycenters, and unbalanced transport~\citep{Cuturi13,benamou2015iterative,chizat2018scaling}. A general convergence blueprint helps identify which properties make these methods effective and which constants deteriorate when the regularization is reduced.

\paragraph{Bregman projections and iterative scaling.}
Alternating projections originate in Bregman's method~\citep{Bregman67} and the information-geometric analysis of Csisz\'ar and Tusn\'ady~\citep{CsiszarTusnady84}; a general convex-analytic treatment is given in~\citep{BauschkeBorwein1997}. Iterative proportional fitting and matrix scaling have a long history~\citep{DemingStephanIPFP,Sinkhorn64}. For positive transport kernels, projective-metric contraction gives linear convergence~\citep{FranklinLorenz89,borwein1994dad}, but the resulting contraction factor can approach one exponentially fast as $\gamma$ decreases. We instead study a weaker quantity, the dual objective gap, with explicit polynomial dependence on the regularization. This does not imply equally sharp primal-distance convergence.

\paragraph{Robust transport bounds.}
Polynomial complexity estimates for scaling and entropic transport were developed in~\citep{kalantari2008complexity,altschuler2017near,dvurechensky2018computational,chakrabarty2021sinkhornsublinearrate}; related perspectives use mirror descent~\citep{leger2021gradient,aubin2022mirror}. For balanced probability marginals, bounded mass leads to an $O(1/(\gamma k))$ dual rate. More generally, a rate written as $O(X_\gamma U_\gamma^2/(\gamma k))$ is only as robust as the bounds on $X_\gamma$ and $U_\gamma$. In particular, our graph-flow bound has $X_\gamma=O(1/\gamma)$, giving $O(1/(\gamma^2k))$ for fixed graph data. Linear-rate refinements under additional assumptions are complementary~\citep{chizat2025sharper}; neither kind of estimate eliminates the statistical and computational tradeoffs of choosing $\gamma$.

\paragraph{Transport on graphs.}
For a shortest-path ground metric, Wasserstein--1 transport admits a minimum-cost flow formulation, the discrete counterpart of Beckmann's formulation~\citep{Beckmann52,AhujaEtAl93}. Tree distances even admit closed-form transport costs~\citep{evans2012phylogenetic}. General networks admit sophisticated exact and approximate flow algorithms~\citep{orlin1993minimum,daitch2008faster}; our contribution is not a better worst-case dependence on accuracy than these solvers. It is a simple sparse entropy-regularized iteration with an explicit convergence analysis. Applying vanilla Sinkhorn to graph transport instead acts on the generally dense shortest-path distance matrix: setting non-edge costs to $+\infty$ would forbid direct coupling of nonadjacent vertices, not allow routing through them, and defines a different transport problem. Dense rounding results~\citep{altschuler2017near} remain applicable to the dense formulation, without removing this quadratic storage and update bottleneck.

\paragraph{Contributions.}
The core contribution is a general blueprint for iterative KL projections and its instantiation in a sparse flow-Sinkhorn algorithm. Figure~\ref{fig:blueprint} links the two main results, Theorems~\ref{thm:kl-dual-rate} and~\ref{thm:approx-linprog}, to the primal and dual helpers, Propositions~\ref{prop:mass-bound-block} and~\ref{prop:uniform-iter-final}. The graph application gives $O(p)$ sweeps and a fully specified optimal-value guarantee (Theorem~\ref{thm:graphw1-complexity}); CPU and GPU experiments on synthetic and genomic-inspired graphs illustrate the tradeoff between regularization bias and runtime. Balanced transport provides a second complete instantiation (Appendix~\ref{app-app:sinkhorn-ot}). Quantum transport is a promising further example, particularly for quantum-information and quantum-machine-learning problems: its trace-norm Pinsker geometry fits the Bregman blueprint, but our explicit quotient-norm/non-expansiveness route is not yet established for its dual block maps. Existing quantum Sinkhorn convergence and dual estimates must be distinguished from this remaining step (Appendix~\ref{app-app:quantum-ot}). Code, benchmarks, and the Lean development are available at \url{https://github.com/gpeyre/flow-sinkhorn}; Appendix~\ref{app-sec:lean-formalization} states the formalization's current limitations. Auxiliary non-expansiveness and general Bregman arguments are in Appendices~\ref{app-sec:non-expansiveness} and~\ref{app-app:general-bregman}.

\begin{figure}[htbp]
\centering
\resizebox{\linewidth}{!}{\begin{tikzpicture}[x=1mm,y=-1mm,font=\scriptsize,
 hypothesis/.style={draw=blue!45!black,fill=blue!7,rounded corners=1mm,align=center,inner sep=2mm,text width=35mm},
 result/.style={draw=red!50!black,fill=red!9,rounded corners=1mm,align=center,inner sep=2mm,text width=35mm},
 helper/.style={draw=black!45,fill=yellow!16,rounded corners=1mm,align=center,inner sep=2mm,text width=35mm},
 graph/.style={draw=green!40!black,fill=green!10,rounded corners=1mm,align=center,inner sep=2mm,text width=35mm},
 barr/.style={-{Stealth[length=1.8mm]},blue!65!black,line width=.65pt},
 rarr/.style={-{Stealth[length=1.8mm]},red!65!black,line width=.65pt}]
\node[font=\bfseries] at (20,0) {Graph inputs};
\node[font=\bfseries] at (67,0) {Helper results};
\node[font=\bfseries] at (114,0) {Blueprint hypotheses};
\node[font=\bfseries] at (161,0) {Conclusions};
\node[graph] (order) at (20,19) {Scalar graph block maps\\Prop.~\ref{prop:graphw1-psi2-closed-nonexp}\\order + constant shifts};
\node[helper,fill=orange!14] (topical) at (67,19) {Prop.~\ref{app-prop:topical-nonexpansive}\\order + shifts $\Rightarrow$\\variation non-expansiveness};
\node[hypothesis] (ne) at (114,19) {$\Psi$ non-expansive\\in the first block quotient};
\node[graph] (h) at (20,48) {Positive cost + reverse edges\\Prop.~\ref{app-prop:hgamma-graphw1}\\$H_\gamma$ explicit};
\node[hypothesis] (hk) at (67,48) {$\|\log(x^\star/z)\|_\infty\leq H_\gamma$\\$\kappa_1$ finite, def.~\eqref{eq:kappa}};
\node[helper,fill=cyan!14] (dual) at (114,48) {Dual helper\\Prop.~\ref{prop:uniform-iter-final}\\$U_\gamma=2\kappa_1(\|C\|_\infty+\gamma H_\gamma)$};
\node[hypothesis] (radius) at (161,48) {First-block orbit and optimum\\quotient radius $\leq U_\gamma$};
\node[graph] (geo) at (20,77) {Paths: $\kappa_1\leq\Dhop$\\Prop.~\ref{app-prop:kappa-graph-diameter}\\$b_2=0$, $\|A\|_{1\to1}=2$};
\node[hypothesis] (pairing) at (67,77) {Pairing bound $\langle b,u\rangle\leq P_\gamma$\\not independent gauges\\graph: $P_\gamma=\|q\|_1U_\gamma$};
\node[helper] (mass) at (114,77) {Primal helper\\Prop.~\ref{prop:mass-bound-block}\\$X_\gamma=P_\gamma/\gamma+\sum_i z_i e^{-C_i/\gamma}$};
\node[hypothesis] (x) at (161,77) {Full and half-step\\masses $\leq X_\gamma$};
\node[graph] (scale) at (20,109) {Normalized graph lift\\cost $(W/2,W/2)$\\parameter $\theta=\gamma/2$};
\node[helper,fill=green!10] (pinsker) at (67,109) {Lemma~\ref{app-lem:pinsker-nonnormalized}\\bounded-mass Pinsker\\$\eta_\gamma=1/(2X_\gamma)$};
\node[result] (rate) at (161,109) {Thm.~\ref{thm:kl-dual-rate}\\$\Delta_k\leq8X_\gamma U_\gamma^2a^2/(\gamma k)$\\$a=\|A\|_{1\to1}$};
\node[hypothesis] (bias) at (114,137) {Bias envelope $B_0$\\Lemma~\ref{app-lem:kl-bias}\\$0\leq F_\gamma^\star-F_0^\star\leq\gamma B_0$};
\node[result] (accuracy) at (161,137) {Thm.~\ref{thm:approx-linprog}\\$\gamma=\varepsilon/(2B_0)$\\$K\geq32X_\gamma U_\gamma^2a^2 B_0/\varepsilon^2$};
\node[graph] (final) at (67,137) {Graph value estimate\\Thm.~\ref{thm:graphw1-complexity}\\fixed data: $O(p\varepsilon^{-3})$};
\draw[barr] (order.east)--(topical.west);
\draw[rarr] (topical.east)--(ne.west);
\draw[barr] (ne.south)--(dual.north);
\draw[rarr] (h.east)--(hk.west);
\draw[rarr] (geo.north east)--(hk.south west);
\draw[barr] (hk.east)--(dual.west);
\draw[rarr] (dual.east)--(radius.west);
\draw[barr] (pairing.east)--(mass.west);
\draw[rarr] (mass.east)--(x.west);
\draw[barr] (x.south)--(rate.north);
\draw[barr] (radius.east)--++(4,0)|-(rate.east);
\draw[barr] (x.south west)--++(-3,7)-|(pinsker.north);
\draw[rarr] (pinsker.east)--(rate.west);
\draw[barr] (geo.south)--(scale.north);
\draw[barr] (rate.south)--(accuracy.north);
\draw[barr] (bias.east)--(accuracy.west);
\draw[rarr] (accuracy.south)--++(0,8)-|(final.south);
\draw[barr] (scale.south)|-(final.west);
\end{tikzpicture}
}
\caption{The convergence blueprint. Blue arrows supply hypotheses to results; red arrows carry their conclusions. Green boxes instantiate the abstract inputs for graph flow. Every bound is stated for the regularization used in the corresponding formulation; for graph flow the lifted parameter is $\theta=\gamma/2$.}
\label{fig:blueprint}
\end{figure}
